\chapter{Referencial Teórico}
\label{cap:referencial}

\section{Considerações Iniciais}

Este capítulo apresenta os conceitos que fundamentam a comparação proposta neste trabalho.
Inicialmente, discute-se o que é arquitetura de software e por que ela é determinante para a
qualidade de um sistema. Em seguida, são descritas as arquiteturas monolítica e de
microsserviços, suas características, vantagens e limitações, bem como as formas de
comunicação entre os componentes de cada uma. Depois, são definidos os atributos de qualidade
utilizados como critérios de comparação e as métricas associadas a eles. Os trabalhos que
compararam as duas abordagens são apresentados separadamente, no
Capítulo~\ref{cap:estado-arte}.

\section{Arquitetura de Software}

Segundo \citeonline{bass2021}, a arquitetura de software de um sistema é o conjunto de
estruturas necessárias para raciocinar sobre ele, compreendendo os elementos de software, as
relações entre eles e as propriedades de ambos. A arquitetura, portanto, não se limita à
organização do código-fonte: ela abrange também a forma como os componentes são executados,
como se comunicam e como são implantados.

Uma característica central da arquitetura é que suas decisões são difíceis e caras de reverter
depois que o sistema está em produção. Por isso, \citeonline{richards2024} defendem que o
trabalho do arquiteto consiste em analisar \textit{trade-offs}: toda decisão arquitetural traz
benefícios e custos, e não existe uma solução universalmente melhor, apenas a mais adequada
para um conjunto específico de requisitos. Esses autores distinguem ainda as
\textit{características arquiteturais}, como desempenho, escalabilidade e manutenibilidade,
dos requisitos funcionais do sistema; são justamente essas características que costumam
orientar a escolha de um estilo arquitetural.

A divisão de um sistema em partes menores é um tema clássico da engenharia de software.
\citeonline{parnas1972} mostrou que a forma de decompor um sistema em módulos afeta
diretamente sua flexibilidade e compreensibilidade, propondo que cada módulo esconda de
outros as decisões de projeto mais sujeitas a mudanças (ocultação de informação). Esse
princípio está na base tanto dos monólitos modulares quanto dos microsserviços. Já
\citeonline{conway1968} observou que a estrutura de um sistema tende a refletir a estrutura de
comunicação da organização que o desenvolve, observação conhecida como Lei de Conway e
frequentemente citada na discussão sobre microsserviços, que costumam ser organizados em torno
de equipes autônomas.

\citeonline{richards2024} classificam os estilos arquiteturais em dois grandes grupos:
\textit{monolíticos}, nos quais o sistema é implantado como uma única unidade, e
\textit{distribuídos}, nos quais o sistema é formado por várias unidades implantadas
separadamente que se comunicam pela rede. As arquiteturas comparadas neste trabalho
representam esses dois grupos e são descritas nas seções seguintes.

\section{Arquitetura Monolítica}

Na arquitetura monolítica, todas as funcionalidades do sistema são construídas, empacotadas e
implantadas como uma única aplicação, executada em um único processo
\cite{richardson2018}. Em uma aplicação Web típica, isso significa que a interface de
programação, as regras de negócio e o acesso a dados residem no mesmo artefato e,
normalmente, compartilham um único banco de dados. A comunicação entre as partes do sistema
ocorre por chamadas de método dentro do próprio processo, sem passar pela rede.

\citeonline{newman2021} distingue três formas de monólito: o \textit{monólito de processo
único}, em que todo o código é implantado como um único processo; o \textit{monólito
modular}, variação do anterior em que o código é dividido em módulos bem delimitados que,
ainda assim, são implantados em conjunto; e o \textit{monólito distribuído}, um sistema
formado por vários serviços que, por estarem fortemente acoplados, precisam ser implantados
juntos, reunindo as desvantagens das duas abordagens. O monólito modular é apontado pelo autor
como uma alternativa sólida, pois preserva a simplicidade operacional do monólito e, ao mesmo
tempo, organiza o código de forma semelhante à decomposição em serviços.

Entre as vantagens da arquitetura monolítica, \citeonline{richardson2018} destaca a
simplicidade de desenvolvimento, de testes e de implantação, já que existe apenas um artefato
a ser construído, testado e publicado. As transações também são mais simples: como os dados
estão em um único banco, operações que envolvem várias partes do sistema podem ser executadas
em uma única transação, com as garantias de atomicidade e consistência oferecidas pelo
sistema gerenciador de banco de dados. Além disso, a ausência de comunicação pela rede entre
as partes do sistema elimina a latência e as falhas parciais típicas dos sistemas
distribuídos.

Por outro lado, à medida que a aplicação cresce, o monólito tende a apresentar limitações.
\citeonline{richardson2018} descreve esse cenário como o ``inferno do monólito'': o código se
torna grande e difícil de compreender, o tempo de compilação e de inicialização aumenta,
qualquer alteração exige a reimplantação de toda a aplicação e a escalabilidade fica restrita
à replicação do sistema inteiro, mesmo quando apenas uma parte dele está sobrecarregada. A
aplicação também fica presa a uma única pilha tecnológica, o que dificulta a adoção de novas
tecnologias.

Apesar dessas limitações, \citeonline{fowler2015} argumenta que quase todos os casos de
sucesso com microsserviços começaram como monólitos que cresceram demais e foram
decompostos, enquanto sistemas construídos como microsserviços desde o início com frequência
enfrentaram sérios problemas. Por isso, o autor recomenda começar com um monólito e só
decompô-lo quando os limites entre as partes do sistema estiverem bem compreendidos.

\section{Arquitetura de Microsserviços}

O termo microsserviços foi popularizado por \citeonline{lewis2014}, que descrevem o estilo como
uma abordagem para desenvolver uma aplicação como um conjunto de pequenos serviços, cada um
executado em seu próprio processo e comunicando-se por mecanismos leves, geralmente uma
interface HTTP. Esses serviços são construídos em torno de capacidades de negócio e podem ser
implantados de forma independente por meio de processos automatizados, com o mínimo de
gerenciamento centralizado, podendo inclusive utilizar linguagens de programação e tecnologias
de armazenamento de dados diferentes.

Os mesmos autores apontam um conjunto de características comuns às arquiteturas de
microsserviços, entre elas: a componentização por meio de serviços; a organização em torno de
capacidades de negócio; a preferência por ``\textit{endpoints} inteligentes e canais simples'',
em que a lógica fica nos serviços e não na infraestrutura de comunicação; a governança e o
gerenciamento de dados descentralizados; a automação da infraestrutura; e o projeto voltado à
tolerância a falhas \cite{lewis2014}. \citeonline{dragoni2017} situam os microsserviços como
uma evolução da arquitetura orientada a serviços, caracterizada por serviços menores, mais
coesos e implantados de forma independente.

Para \citeonline{newman2021}, a característica essencial de um microsserviço é poder ser
implantado de forma independente, sem exigir a implantação de outros serviços. Para isso, cada
serviço deve ser modelado em torno de um domínio de negócio e encapsular seus próprios dados.
A definição dos limites de cada serviço costuma se apoiar no conceito de contexto delimitado
(\textit{bounded context}), proposto por \citeonline{evans2003} no projeto orientado a domínio
(\textit{Domain-Driven Design}). Desse princípio decorre o padrão de banco de dados por
serviço, no qual nenhum serviço acessa diretamente os dados de outro; o compartilhamento de
informações ocorre apenas por meio das interfaces expostas por cada serviço
\cite{richardson2018}.

Como os clientes da aplicação precisariam conhecer o endereço de vários serviços, é comum o
uso de um \textit{API Gateway}: um ponto de entrada único que recebe as requisições externas e
as encaminha ao serviço responsável, podendo concentrar responsabilidades transversais, como a
autenticação \cite{richardson2018}. A Figura~\ref{fig:arquiteturas} ilustra a diferença
estrutural entre as duas arquiteturas, utilizando como exemplo os seis módulos da aplicação de
\textit{e-commerce} desenvolvida neste trabalho: autenticação (Auth), participantes, produtos,
loja, pedidos e pagamentos. No monólito, esses módulos compartilham o mesmo processo e o mesmo
banco de dados; nos microsserviços, cada um se torna um serviço com banco próprio, acessado
pelo \textit{API Gateway} e integrado aos demais por requisições HTTP ou por mensagens
publicadas em um \textit{broker}.

\begin{figure}[htbp]
\centering
\caption{Organização de uma aplicação de \textit{e-commerce} nas arquiteturas (a) monolítica e (b) de microsserviços}
\label{fig:arquiteturas}
\resizebox{\textwidth}{!}{%
\begin{tikzpicture}[
    modulo/.style={draw, rounded corners=2pt, minimum width=2.6cm, minimum height=0.8cm, fill=gray!10, font=\small},
    bd/.style={draw, cylinder, shape border rotate=90, aspect=0.25, minimum width=1.4cm, minimum height=0.9cm, fill=gray!25, font=\footnotesize},
    seta/.style={-{Stealth[length=2.2mm]}, thick},
    rotulo/.style={font=\small\bfseries, anchor=west}
]
\node[rotulo] at (-1.3, 2.7) {(a) Arquitetura monolítica};
\node[modulo] (mauth) at (0,0) {Auth};
\node[modulo] (mpart) at (2.9,0) {Participante};
\node[modulo] (mprod) at (5.8,0) {Produto};
\node[modulo] (mloja) at (8.7,0) {Loja};
\node[modulo] (mped) at (11.6,0) {Pedido};
\node[modulo] (mpag) at (14.5,0) {Pagamento};
\node[draw, thick, rounded corners=4pt, inner sep=8pt, fit=(mauth)(mpag), label={[font=\footnotesize]below:Aplicação única: um processo e um artefato de implantação}] (mono) {};
\node[modulo, fill=white] (mcli) at (7.25,1.7) {Cliente};
\draw[seta] (mcli) -- (mono.north -| mcli);
\node[bd, minimum width=2.2cm] (mbd) at (7.25,-2.5) {BD único};
\draw[seta] (mono.south -| mbd) ++(0,-0.45) -- (mbd.north);

\node[rotulo] at (-1.3, -3.8) {(b) Arquitetura de microsserviços};
\node[modulo, fill=white] (cli) at (7.25,-4.5) {Cliente};
\node[modulo, fill=gray!30] (gw) at (7.25,-5.9) {API Gateway};
\draw[seta] (cli) -- (gw);
\foreach \nome/\x/\id in {Auth/0/sauth, Participante/2.9/spart, Produto/5.8/sprod, Loja/8.7/sloja, Pedido/11.6/sped, Pagamento/14.5/spag} {
    \node[modulo] (\id) at (\x,-7.6) {\nome};
    \node[bd] (bd\id) at (\x,-9.2) {BD};
    \draw[seta] (\id) -- (bd\id);
    \draw[seta] (gw.south) -- (\id.north);
}
\node[draw, rounded corners=2pt, minimum width=17.1cm, minimum height=0.8cm, fill=gray!30, font=\small] (broker) at (7.25,-10.9) {\textit{Broker} de mensagens (comunicação assíncrona)};
\foreach \id in {sauth, spart, sprod, sped, spag} {
    \draw[dashed, thick, {Stealth[length=2mm]}-{Stealth[length=2mm]}] ([xshift=-0.15cm]\id.south east) -- ([xshift=-0.15cm]\id.south east |- broker.north);
}
\end{tikzpicture}%
}
\par\vspace{4pt}{\footnotesize Fonte: Elaborada pelo autor (2026).}
\end{figure}

Entre os benefícios dos microsserviços, \citeonline{newman2021} destaca a possibilidade de
utilizar tecnologias diferentes em cada serviço, a maior robustez (uma falha em um serviço
não precisa derrubar todo o sistema), a escalabilidade seletiva (apenas os serviços
sobrecarregados precisam ser replicados), a facilidade de implantar alterações de forma
isolada e o alinhamento entre a arquitetura e a organização das equipes.

Esses benefícios, entretanto, têm custos. Ao distribuir o sistema, a aplicação passa a lidar
com problemas que não existem em um monólito: a comunicação pela rede é mais lenta e sujeita a
falhas; os dados ficam espalhados entre vários bancos, o que impede transações locais e
geralmente leva à consistência eventual; e a quantidade de componentes a serem configurados,
monitorados, testados e implantados aumenta consideravelmente \cite{newman2021}. Para manter a
consistência em operações que envolvem vários serviços, \citeonline{richardson2018} propõe o
padrão \textit{Saga}, em que a operação é dividida em uma sequência de transações locais
coordenadas por mensagens. \citeonline{richards2024} lembram ainda as chamadas falácias da
computação distribuída, suposições equivocadas como a de que a rede é confiável, a latência
é zero e a largura de banda é infinita, que, quando ignoradas, comprometem sistemas
distribuídos como os baseados em microsserviços. A Tabela~\ref{tab:comparacao-arquiteturas}
sintetiza as principais diferenças entre as duas arquiteturas.

\begin{table}[htbp]
\centering
\caption{Síntese das principais diferenças entre as arquiteturas monolítica e de microsserviços}
\label{tab:comparacao-arquiteturas}
\small
\begin{tabular}{|>{\raggedright\arraybackslash}p{3.2cm}|>{\raggedright\arraybackslash}p{5.9cm}|>{\raggedright\arraybackslash}p{5.9cm}|}
\hline
\textbf{Aspecto} & \textbf{Monolítica} & \textbf{Microsserviços} \\ \hline
Implantação & Um único artefato; toda alteração exige reimplantar o sistema inteiro & Cada serviço é implantado de forma independente \\ \hline
Escalabilidade & Replicação da aplicação inteira & Replicação apenas dos serviços necessários \\ \hline
Comunicação interna & Chamadas de método no mesmo processo & Chamadas pela rede (HTTP) e mensageria \\ \hline
Dados & Banco de dados compartilhado & Banco de dados por serviço \\ \hline
Consistência & Transações locais (ACID) & Consistência eventual e padrões como \textit{Saga} \\ \hline
Falhas & Uma falha grave pode indisponibilizar todo o sistema & Falhas podem ficar isoladas em um serviço, mas surgem falhas parciais de rede \\ \hline
Tecnologia & Pilha tecnológica única & Possibilidade de tecnologias diferentes por serviço \\ \hline
Complexidade operacional & Menor: poucos componentes a implantar e monitorar & Maior: muitos componentes, infraestrutura de comunicação e observabilidade \\ \hline
\end{tabular}
\par\vspace{4pt}{\footnotesize Fonte: Elaborada pelo autor com base em \citeonline{lewis2014}, \citeonline{richardson2018} e \citeonline{newman2021}.}
\end{table}

\section{Comunicação entre Componentes}

A forma como as partes de um sistema se comunicam é uma das principais diferenças entre as
duas arquiteturas. Em um monólito, os módulos se comunicam por chamadas de método dentro do
mesmo processo, que são rápidas, confiáveis e podem participar da mesma transação. Em uma
arquitetura de microsserviços, os serviços precisam se comunicar pela rede, de forma síncrona
ou assíncrona \cite{richardson2018}.

\subsection{Comunicação Síncrona}

Na comunicação síncrona, um serviço envia uma requisição e aguarda a resposta para continuar o
processamento. A forma mais comum desse tipo de comunicação em microsserviços é o uso de
requisições HTTP seguindo o estilo arquitetural REST (\textit{Representational State
Transfer}), definido por \citeonline{fielding2000} como um conjunto de restrições para sistemas
hipermídia distribuídos, entre elas a comunicação cliente-servidor sem estado e o uso de uma
interface uniforme para manipular recursos.

A comunicação síncrona é simples de implementar e de compreender, pois se assemelha a uma
chamada de função. No entanto, ela cria um acoplamento temporal entre os serviços: o serviço
que faz a chamada depende de que o outro esteja disponível e responda a tempo. Assim, a
latência de cada chamada se soma ao tempo total da operação e a indisponibilidade de um
serviço pode se propagar em cascata para os serviços que dependem dele
\cite{newman2021, richardson2018}.

\subsection{Comunicação Assíncrona e Mensageria}

Na comunicação assíncrona, o serviço que origina a comunicação não aguarda uma resposta
imediata. Em vez disso, publica uma mensagem em um intermediário, o \textit{broker} de
mensagens, que a entrega aos serviços interessados. \citeonline{hohpe2003} descrevem a
mensageria como uma forma de integração que desacopla emissores e receptores: eles não
precisam conhecer um ao outro nem estar disponíveis ao mesmo tempo, pois as mensagens ficam
armazenadas no intermediário até serem consumidas.

O Apache Kafka, utilizado neste trabalho, é uma plataforma de mensageria distribuída
originalmente desenvolvida no LinkedIn para o processamento de grandes volumes de dados de
\textit{log} \cite{kreps2011}. Nele, as mensagens são publicadas em tópicos e mantidas em um
\textit{log} persistente, a partir do qual os consumidores as leem no próprio ritmo. Esse
modelo permite que um serviço reaja a eventos de outro, como a criação de um usuário ou a
aprovação de um pagamento, sem que exista uma dependência direta entre eles.

Em contrapartida, a comunicação assíncrona torna o fluxo do sistema mais difícil de
acompanhar e depurar, introduz um componente adicional de infraestrutura e leva à
consistência eventual: durante um intervalo de tempo, os dados de serviços diferentes podem
não refletir o mesmo estado \cite{richardson2018}.

\section{Atributos de Qualidade Avaliados}

A qualidade de um produto de software pode ser descrita por meio de atributos não funcionais.
A norma ISO/IEC 25010 define um modelo de qualidade de produto composto, em sua edição de 2023,
por nove características, entre elas a eficiência de desempenho, a manutenibilidade e a
flexibilidade, esta última incluindo a escalabilidade como subcaracterística
\cite{iso25010}. \citeonline{bass2021} tratam esses atributos como \textit{atributos de
qualidade} e mostram que cada decisão arquitetural favorece alguns deles em detrimento de
outros. Esta seção descreve os cinco atributos utilizados como critérios de comparação neste
trabalho e as métricas associadas a cada um.

\subsection{Desempenho}

O desempenho diz respeito à capacidade do sistema de responder às requisições dentro dos
limites de tempo esperados e de utilizar os recursos computacionais de forma eficiente. A
norma ISO/IEC 25010 trata esse atributo como \textit{eficiência de desempenho}, dividida em
comportamento temporal, utilização de recursos e capacidade \cite{iso25010}.

\citeonline{jain1991} apresenta as métricas mais utilizadas na análise de desempenho de
sistemas computacionais: o \textit{tempo de resposta}, intervalo entre o envio de uma
requisição e o recebimento da resposta; a \textit{vazão} (\textit{throughput}), quantidade de
requisições atendidas por unidade de tempo; e a \textit{utilização}, fração do tempo em que um
recurso, como o processador, permanece ocupado. Como o tempo de resposta varia entre
requisições, é comum analisá-lo por meio de percentis: o percentil 95, por exemplo, indica o
tempo dentro do qual 95\% das requisições foram atendidas, revelando o comportamento das
requisições mais lentas, que a média tende a ocultar.

Na comparação entre as arquiteturas, espera-se que a comunicação pela rede entre os
microsserviços acrescente latência às operações que envolvem vários serviços, enquanto no
monólito essas operações ocorrem em memória \cite{richards2024}.

\subsection{Escalabilidade}

\citeonline{bondi2000} define escalabilidade como a capacidade de um sistema de acomodar o
aumento de carga, de volume de dados ou de usuários mantendo um desempenho aceitável. Há duas
formas principais de escalar um sistema: a \textit{escalabilidade vertical}, que aumenta os
recursos de uma mesma máquina, como processador e memória; e a \textit{escalabilidade
horizontal}, que adiciona novas instâncias do sistema e distribui a carga entre elas.

\citeonline{abbott2015} propõem o cubo de escalabilidade, que organiza as estratégias de
escala em três eixos: no eixo X, a aplicação inteira é replicada atrás de um balanceador de
carga; no eixo Y, o sistema é decomposto por funcionalidade, de modo que cada parte possa ser
escalada separadamente; e no eixo Z, os dados são particionados entre instâncias.
\citeonline{richardson2018} utiliza esse modelo para mostrar que o monólito escala
essencialmente no eixo X, enquanto os microsserviços correspondem a uma decomposição no eixo
Y, permitindo replicar apenas os serviços sob maior demanda. A escalabilidade pode ser medida
pela variação da vazão e do tempo de resposta à medida que se acrescentam instâncias e pela
eficiência dessa escala, isto é, o ganho obtido em relação aos recursos adicionados.

\subsection{Manutenibilidade}

A manutenibilidade é o grau de eficácia e eficiência com que um sistema pode ser modificado
para corrigir defeitos, melhorá-lo ou adaptá-lo. A norma ISO/IEC 25010 a divide em
modularidade, reusabilidade, analisabilidade, modificabilidade e testabilidade
\cite{iso25010}.

Para quantificar esse atributo, são utilizadas métricas de código-fonte. A
\textit{complexidade ciclomática}, proposta por \citeonline{mccabe1976}, mede o número de
caminhos independentes no fluxo de controle de um programa: quanto maior, mais difícil é
compreender e testar o código. \citeonline{chidamber1994} propuseram um conjunto de métricas
para sistemas orientados a objetos, conhecido como métricas CK, entre as quais se destacam o
acoplamento entre objetos (CBO), que conta de quantas outras classes uma classe depende, e a
falta de coesão em métodos (LCOM), que indica se os métodos de uma classe trabalham sobre os
mesmos dados. Somam-se a essas métricas o tamanho do código, medido em linhas de código, e a
duplicação. Uma forma complementar de avaliar a manutenibilidade é analisar o impacto de uma
mesma alteração nas duas arquiteturas, verificando quantos arquivos, módulos e implantações
são afetados.

\subsection{Complexidade de Desenvolvimento}

A complexidade de desenvolvimento refere-se ao esforço necessário para construir e evoluir o
sistema, considerando a quantidade de componentes, configurações e conceitos técnicos que os
desenvolvedores precisam dominar. Diferentemente dos demais atributos, não se trata de uma
característica definida na norma ISO/IEC 25010, mas de um aspecto recorrente na comparação
entre as duas arquiteturas.

\citeonline{newman2021} aponta que os microsserviços afetam a experiência dos desenvolvedores
e podem levar a uma sobrecarga tecnológica, pois exigem o domínio de ferramentas de
comunicação, mensageria, implantação e observabilidade que não seriam necessárias em um
monólito. Parte dessa complexidade não decorre do problema de negócio em si, mas da
distribuição do sistema, como o tratamento de falhas de rede, a autenticação entre serviços e
a duplicação de modelos de dados entre serviços. Essa complexidade pode ser estimada pelo
número de projetos e artefatos, pela quantidade de arquivos de configuração, pelo número de
dependências e pela proporção de código dedicada à infraestrutura de comunicação.

\subsection{Implantação}

A implantação (\textit{deployability}) é a capacidade de colocar uma nova versão do sistema em
execução de forma rápida, previsível e com baixo risco \cite{bass2021}. Nos microsserviços,
cada serviço pode ser implantado de forma independente, o que reduz o impacto de cada
alteração, mas multiplica o número de artefatos, processos e componentes de infraestrutura que
precisam ser gerenciados \cite{newman2021}. No monólito, existe um único artefato, mais simples
de implantar, porém qualquer alteração exige a reimplantação de todo o sistema.

Esse atributo pode ser avaliado pelo número de artefatos e contêineres necessários para
executar o sistema, pelo tempo de construção e de inicialização, pelo tamanho das imagens de
implantação, pelo consumo de recursos em repouso e pela quantidade de etapas necessárias para
colocar o ambiente em funcionamento.

\section{Considerações Finais}

Este capítulo apresentou os fundamentos teóricos necessários para a comparação proposta. A
arquitetura monolítica oferece simplicidade de desenvolvimento, de implantação e de
consistência de dados, mas tende a se tornar mais difícil de manter e escalar à medida que
cresce. A arquitetura de microsserviços favorece a implantação independente e a escalabilidade
seletiva, ao custo de introduzir a complexidade dos sistemas distribuídos. Os atributos de
qualidade e as métricas descritos servem de base para os critérios de comparação definidos no
Capítulo~\ref{cap:metodologia}. O próximo capítulo apresenta o estado da arte, reunindo os
trabalhos que já compararam as duas arquiteturas na prática.
